\section{Period and Cycle Reconstruction}
\label{sec:period-and-cycle-reconstruction}

The finite representation works because the filter predicate is periodic
and the scan preserves the order of the repeated accepted residues.

\begin{itemize}
  \item one period shifts every generated value by $M$;
  \item every later block is a translated copy of the first block;
  \item integrating one finite positive gap cycle reconstructs the
    infinite scan.
\end{itemize}

\subsection{Canonical Block Shift}
\label{sec:canonical-block-shift}

The boundary $h+M$ passes exactly the same filters as $h$, so
completeness gives a positive index $T$ with $\ell_T=h+M$. Periodicity of
acceptance and strict order then identify the $k$-th accepted value in
the following block:

\begin{equation*}
\begin{aligned}
A_S(v+M) &= A_S(v)
  &&\text{[Modulo-period invariance]} \\
\ell_T &= h+M
  &&\text{[Canonical boundary]} \\
\ell_{k+T} &= \ell_k+M
  &&\text{[Same ordered survivor]} \\
\ell_{k+nT} &= \ell_k+nM
  &&\text{[Block induction]}\quad\blacksquare\ \text{[Q.E.D.]}.
\end{aligned}
\end{equation*}

This property is verified in
\href{https://github.com/thiagomata/prime-numbers/blob/master/src/main/scala/v1/chapter6/sieve/seq/spec/properties/SpecSieveSeqPeriodProperties.scala}{\texttt{SpecSieveSeqPeriodProperties::assertBlockShiftMultiple}}.

\subsection{Gap-Cycle Reconstruction}
\label{sec:gap-cycle-reconstruction}

Let \texttt{GapCycle(G)} store the first $T$ adjacent differences. The
block-shift theorem makes those differences periodic:

\begin{equation*}
\begin{aligned}
g_{k+T}
&=\ell_{k+T+1}-\ell_{k+T} \\
&=(\ell_{k+1}+M)-(\ell_k+M)
  &&\text{[Block shift]} \\
&=g_k.
\end{aligned}
\end{equation*}

The cycle integral starts at $h$ and adds these gaps. Induction on $k$
then reconstructs every scan value:

\begin{equation*}
\begin{aligned}
I_G(0)
&=h+g_0=\ell_1
  &&\text{[Base]} \\
I_G(k)
&=I_G(k-1)+g_k \\
&=\ell_k+(\ell_{k+1}-\ell_k)
  &&\text{[Induction hypothesis]} \\
&=\ell_{k+1}
  &&\blacksquare\ \text{[Q.E.D.]}.
\end{aligned}
\end{equation*}

{\raggedright
This property is verified in
\href{https://github.com/thiagomata/prime-numbers/blob/master/src/main/scala/v1/chapter6/sieve/seq/spec/properties/SpecSieveSeqPeriodProperties.scala}{\texttt{SpecSieveSeqPeriodProperties::assertSpecGapCycleIntegralMatchesApply}}.
\par}

\subsection{Repetition Does Not Change the Infinite Sequence}
\label{sec:repetition-invariance}

The transition prepares $h$ copies of the current period before
installing the head $h$ as a new filter. Repeating the stored gap list
multiplies the finite representation's period from $T$ to $hT$, but it
does not change the infinite periodic sequence represented by the cycle
integral.

\begin{equation*}
\begin{aligned}
G^{\langle h\rangle}
  &=\underbrace{G\mathbin{\texttt{++}}\cdots\mathbin{\texttt{++}}G}_{h\text{ copies}}, \\
|G^{\langle h\rangle}|&=hT, \\
G^{\langle h\rangle}_{k\bmod hT}
  &=G_{k\bmod T}, \\
I_{G^{\langle h\rangle}}(k)
  &=I_G(k)
  \quad\text{[Equal increments and initial value]}\quad\blacksquare\ \text{[Q.E.D.]}.
\end{aligned}
\end{equation*}

{\raggedright
This property is verified in
\href{https://github.com/thiagomata/prime-numbers/blob/master/src/main/scala/v1/chapter6/sieve/seq/spec/properties/SpecDerivedRepeatedCycleProperties.scala}{\texttt{SpecDerivedRepeatedCycleProperties::assertSpecRepeatedCycleIntegralMatchesBase}}.
\par}
